%------------------------------------------------------------------------------%
\documentclass[a4paper,12pt,fleqn]{article}

\usepackage{integracao_series}

\ShowAnswers

%------------------------------------------------------------------------------%
\begin{document}

\makeHeader{Integração e Séries}{Prova 1 -- Turma 3}{01/04/2026}

% 01
%------------------------------------------------------------------------------%
\exercise{40\%}
Calcule as integrais
\begin{enumerate}[label=\alph*),itemsep=0pt,topsep=0pt]
  \item\!\! [20\%] \(\int a(3x-1)^4 - \frac{1}{cx} + \sqrt{x} + 7 \, dx\)
  \item\!\! [20\%] \(\int_0^x x\sin(\pi t) \, dt \)
\end{enumerate}

\clearpagequestiononly

\begin{answer}
\textbf{a)} \quad
  Esta é uma integral indefinida, isto é, ela é a função que quando derivada
  retorna o integrando
  \begin{align*}
    F(x)
    & = \int a(3x-1)^4 - \frac{1}{cx} + \sqrt{x} + 7 \, dx
    \\[2mm]
    & = a\int (3x-1)^4 \, dx
      - \frac{1}{c}\int \frac{1}{x} \, dx
      + \int \sqrt{x} \, dx
      + 7 \int dx
    \\[2mm]
    & = a\int (3x-1)^4 \, dx
      - \frac{1}{c}\int x^{-1} \, dx
      + \int x^{\sfrac{1}{2}} \, dx
      + 7\int x^0 \, dx
    \\[2mm]
    & = a\frac{(3x-1)^{4+1}}{3(4+1)}
      - \frac{1}{c}\ln\abs{x}
      + \frac{x^{\sfrac{1}{2}+1}}{\sfrac{1}{2}+1}
      + 7\frac{x^{0+1}}{0+1} + k
      \\[2mm]
    & = \frac{a}{15} (3x-1)^5
      - \frac{1}{c}\ln\abs{x}
      + \frac{x^{\sfrac{3}{2}}}{\sfrac{3}{2}}
      + 7x + k
      \\[2mm]
    & = \frac{a}{15} (3x-1)^5
      - \frac{1}{c}\ln\abs{x}
      + \frac{2x^{\sfrac{3}{2}}}{3}
      + 7x + k
  \end{align*}

\begin{multicols}{2}
\textbf{b)} \quad
  Esta é uma integral definida então ela é um número real.
  Como vamos usar o Teorema Fundamental do Cálculo para avaliar a integral
  precisamos encontrar a primitiva primeiro.

  \textbf{Atenção:} note que a integral é em $t$, portanto $x$ é uma \textbf{constante}!
  \begin{align*}
    F(t)
    & = \int x\sin(\pi t) \, dt \\[2mm]
    & = x \int \sin(\pi t) \, dt \\[2mm]
    & = x \left(\frac{-\cos(\pi t)}{\pi}\right) + c \\[2mm]
    & = -\frac{x}{\pi} \cos(\pi t) + c
  \end{align*}
  Avaliando a integral definida
  \begin{align*}
    A
    & = \int_0^x x\sin(\pi t) \, dt \\[2mm]
    & = F(t) \evalat{0}{x}
    \;=\; F(x) - F(0) \\[2mm]
    & = \left(-\frac{x}{\pi} \cos(\pi t) + c\right) \evalat{0}{x}  \\[2mm]
    & = \left(-\frac{x}{\pi} \cos(\pi x) + c\right)
      - \left(-\frac{x}{\pi} \cos(\pi 0) + c\right)   \\[2mm]
    & = -\frac{x}{\pi} \cos(\pi x)
      +  \frac{x}{\pi}  \\[2mm]
    & = \frac{x}{\pi} \left(1 - \cos(\pi x)\right)
  \end{align*}
\end{multicols}
\end{answer}

% 02
%------------------------------------------------------------------------------%
\exercise{30\%}
Calcule a área entre as curvas
\(
  y = x^2
\)
e
\(
  y = 8\sqrt{x}
\)

\clearpagequestiononly

\begin{answer}
\begin{multicols}{2}
\centerline{\includegraphics{T3_A_area}}

Intersecção entre as curvas
\begin{align*}
  x^2                     & = 8\sqrt{x}                \\
  x^4                     & = \left(8\sqrt{x}\right)^2 \\
  x^4                     & = 8^2 x                    \\
  x^4 - 8^2 x             & = 0                        \\
  x\left(x^3 - 8^2\right) & = 0
\end{align*}
portanto $x=0$ ou
\begin{align*}
  x^3 - 8^2 & = 0                   \\
  x^3       & = 8^2 = 2^{3\times 2} \\
  x         & = 2^2 = 4
\end{align*}
Assim as intersecções ocorrem em $x=0$ e $x=4$

Como no intervalo $[0, 4]$
ambas as funções são contínuas e
$x^2 \leq 8\sqrt{x}$
a área é dada pela integral
\[
  A = \int_0^4 8\sqrt{x} - x^2 \; dx
\]
Calculando a primitiva
\begin{align*}
  F(x)
  & = \int 8\sqrt{x} - x^2 \; dx \\
  & = 8 \int x^{\sfrac{1}{2}} \, dx  - \int x^2 \, dx  \\
  & = 8 \frac{x^{\sfrac{1}{2}+1}}{\sfrac{1}{2}+1}  - \frac{x^{2+1}}{2+1} + c \\
  & = 8 \frac{x^{\sfrac{3}{2}}}{\sfrac{3}{2}}  - \frac{x^3}{3} + c \\
  & = 16 \frac{x^{\sfrac{3}{2}}}{3}  - \frac{x^3}{3} + c \\
  & = \frac{16 x^{\sfrac{3}{2}} - x^3 }{3} + c
\end{align*}
Aplicando o Teorema Fundamental do Cálculo, podemos calcular a área
\begin{align*}
    A
    & = \int_0^4 8\sqrt{x} - x^2 \; dx \\
    & = F(x) \evalat{0}{4} \\
    & F(4) - F(0) \\
    & = \left(\frac{16 \times 4^{\sfrac{3}{2}} - 4^3 }{3}\right) -
        \left(\frac{16 \times 0^{\sfrac{3}{2}} - 0^3 }{3}\right)  \\
    & = \frac{16 \times 8 - 64}{3}  \\
    & = \frac{128 - 64}{3}  \\
    & = \frac{64}{3}
\end{align*}
\end{multicols}
\end{answer}

% 03
%------------------------------------------------------------------------------%
\exercise{30\%}
Considere a região do plano, $R$, entre as curvas
\(
  y = x(x-2)
\)
e
\(
  y = 0
\).
Escreva formulas integrais para os sólidos de revolução construídos:
\begin{enumerate}[label=\alph*),itemsep=0pt,topsep=0pt]
  \item\!\! [15\%] pela rotação da região $R$ em torno do eixo $y$
  \item\!\! [15\%] pela rotação da região $R$ em torno do eixo $x$
\end{enumerate}

\begin{answer}
\begin{multicols}{2}
  \centerline{\includegraphics{T3_A_volume}}

  A intersecção entre as curvas ocorre quando
  \[
    x(x-2) = 0
  \]
  ou seja, quando $x=0$ e $x=2$.
  Assim os intervalos de integração serão
  \[
    a = 0
    \qquad\text{e}\qquad
    b = 2
  \]

  No intervalo $[0, 2]$ as funções são contínuas e
  \[
    x(x-2) \leq 0
  \]
  portanto a ``altura'' entre as curvas é
  \[
    g(x) = 0 - x(x-2) = 2x - x^2
  \]

\textbf{a)} \quad
  Para calcularmos o volume do sólido de revolução criado pela rotação em
  torno do eixo $y$ usamos cascas cilíndricas.
  Neste caso o raio e a altura da casca cilíndrica são
  \[
    r(x) = x
    \qquad
    h(x) = g(x) = 2x - x^2
  \]
  Portanto o volume é
  \begin{align*}
    V
    & = \int_a^b 2\pi r(x) h(x)\, dx \\
    & = \int_0^2 2\pi x\left(2x-x^2\right) \, dx \\
    & = 2\pi \int_0^2 2x^2 - x^3 \, dx
  \end{align*}

\textbf{b)} \quad
  Para calcularmos o volume do sólido de revolução criado pela rotação em
  torno do eixo $x$ usamos seções transversais.
  A área da seção transversão é
  \[
    A(x) = \pi r^2(x) = \pi g^2(x) = \pi\left(2x-x^2\right)^2
  \]
  Assim o volume é dado por
  \begin{align*}
    V
    & = \int_a^b A(x) \, dx \\
    & = \int_0^2 \pi \left(2x-x^2\right)^2 \, dx \\
    & = \pi \int_0^2  \left(2x-x^2\right)^2 \, dx
  \end{align*}
\end{multicols}
\end{answer}

%------------------------------------------------------------------------------%
\end{document}
%------------------------------------------------------------------------------%
